% Indicate the main file. Must go at the beginning of the file.
% !TEX root = ../main.tex

%----------------------------------------------------------------------------------------
% CHAPTER 2
%----------------------------------------------------------------------------------------

\chapter{Methodik}
\label{Chapter2}

%----------------------------------------------------------------------------------------

\section{Projektmanagement und Projektplan}
\label{sec:projektmanagement}

Die Arbeit wurde innerhalb einer Zweiergruppe arbeitsteilig organisiert. Die Bearbeitung war auf vier Wochen
angelegt und orientiert sich an der in der Aufgabenstellung vorgegebenen Gliederung.

\begin{table}[h]
\centering
\caption{Projektplan mit Meilensteinen und Tätigkeiten}
\label{tab:projektplan}
\begin{tabularx}{\textwidth}{p{1cm}p{2.8cm}X}
\toprule
Woche & Meilenstein & Tätigkeiten und Verantwortlichkeiten \\
\midrule
1 & Planung, Daten\-beschaffung &
Sichtung der Aufgabenstellung und des SHP User Guide~\cite{Voorpostel2018}.
Registrierung bei SWISSUbase und Download der Rohdaten.
Festlegung der Ordnerstruktur und Erstellung des Git-Repositories.
Auswahl der zusätzlichen Einflussvariable Zivilstand.
Gemeinsame Arbeit, Koordination durch beide Teammitglieder. \\
\midrule
2 & Daten\-aufbereitung &
Programmierung der Skripte zum Entpacken, Einlesen, Bereinigen, Mergen und Filtern der SHP-Dateien.
Konstruktion der Zielvariablen.
Speicherung des aufbereiteten Datensatzes.
Hauptverantwortung Programmierung: Berdelis; Review und Dokumentation: Schwarz. \\
\midrule
3 & Qualität, Plausibilität, Auswertung &
Analyse fehlender Werte und LOCF-Imputation \footnotemark.
Plausibilitätsvergleich mit BFS-Referenzdaten.
Deskriptive Auswertung der Zielvariable nach Untergruppen.
Hypothesenprüfung zum Zivilstand.
Hauptverantwortung Auswertung: Schwarz; Review und Dokumentation: Berdelis. \\
\midrule
4 & Bericht, \newline Abschluss &
Zusammenführung der Teilergebnisse in diesem Bericht.
Redaktion, Schlusskontrolle, Abgabe.
Gemeinsame Arbeit. \\
\bottomrule
\end{tabularx}
\end{table}
\footnotetext{Imputation ersetzt fehlenden Wert mit letztem früheren verfügbaren Wert.}
Der Programmcode wird mit Git versioniert und über ein gemeinsames Repository\footnote{\href{https://github.com/vvs0x/swiss-household-panel}{https://github.com/vvs0x/swiss-household-panel}} geteilt. Zur Qualitätssicherung
wurde jedes Skript vor dem Zusammenführen in den Hauptzweig vom jeweils anderen
Teammitglied gesichtet und anschliessend verbesserungen vorgenommen.

%----------------------------------------------------------------------------------------
\newpage
\section{Beschreibung der Daten, Beschaffung und Storage}
\label{sec:daten}

\subsection{Datenquelle und Beschaffung}
\label{subsec:datenquelle}

Die Rohdaten stammen aus dem Swiss Household Panel (SHP), das vom Schweizerischen
Kompetenzzentrum für Sozialwissenschaften FORS in Lausanne betrieben wird
\cite{FORS2024}. Der Zugriff erfolgt nach Anmeldung mit der
akademischen edu-ID und Zustimmung zu den Nutzungsbedingungen von FORS, die eine
Verwendung ausschliesslich zu Forschungs- und Lehrzwecken vorsehen.

Von SWISSUbase wurde das vollständige Datenpaket im SPSS-Format heruntergeladen.
Zusätzlich wurde der SHP User Guide in der Fassung bezogen und zur Dokumentation
hinterlegt \cite{Voorpostel2018}. Die Rohdaten werden nicht verändert; sämtliche
Aufbereitungsschritte erzeugen daraus abgeleitete Datensätze in einem separaten
Verzeichnis.

Die für die Plausibilisierung verwendeten Referenzzahlen zur ständigen Wohnbevölkerung
werden direkt über die API des Bundesamts für Statistik im PX-Format eingelesen
\cite{BFS2023}. Auf diese Weise ist jederzeit nachvollziehbar, welcher Stand der
BFS-Daten in den Vergleich eingeflossen ist.

\subsection{Ordnerstruktur und Storage}
\label{subsec:ordnerstruktur}

Das Projektverzeichnis ist funktional gegliedert und trennt Rohdaten, aufbereitete
Daten, Programmcode und Dokumentation. Tabelle~\ref{tab:ordnerstruktur} gibt einen
Überblick über die verwendete Struktur.

\begin{table}[h]
\centering
\caption{Ordnerstruktur des Projektrepositories}
\label{tab:ordnerstruktur}
\begin{tabularx}{\textwidth}{lX}
\toprule
Verzeichnis & Inhalt \\
\midrule
data/raw/           & Unveränderte SHP-Rohdaten im SPSS-Format. \\

data/processed/     & Aus der Aufbereitung erzeugte Zwischen- und Enddatensätze im CSV- und RDS-Format. \\

src/                & R-Markdown-Skripte für Datenaufbereitung und Auswertung sowie die aus ihnen erzeugten PDF-Berichte. \\

docs/bericht/       & Dokumentation: Aufgabenstellung, SHP User Guide sowie das \LaTeX-Projekt für den vorliegenden technischen Bericht. \\

figures/            & Ausgabeverzeichnis für Grafiken, sofern diese ausserhalb der R-Markdown-Berichte verwendet werden. \\

README.md           & Kurzbeschreibung des Repositories, Hinweise zur Reproduktion der Aufbereitung. \\

.gitignore          & Alle Rohdaten sind im .gitignore aufgrund der Forderungen von FORS. \\
\bottomrule
\end{tabularx}
\end{table}

Grosse Binärdateien (insbesondere die Rohdaten und abgeleitete RDS-Dateien) werden über 
von der Versionierung ausgeschlossen, um das Repository schlank zu halten und
urheberrechtliche Anforderungen von FORS einzuhalten. Die Rohdaten verbleiben
ausschliesslich lokal. Alle Ergebnisse lassen sich mit den im Repository hinterlegten
R-Skripten aus der ursprünglichen ZIP-Datei reproduzieren.

\subsection{Verwendete Dateien}
\label{subsec:dateien}

Für die Aufgabe werden vier Dateien des SHP-Datensatzes verwendet. Es handelt sich um
zwei Long-Files, die Beobachtungen über mehrere Wellen hinweg in einer gestapelten
Form enthalten, sowie um zwei Master-Files mit zeitinvarianten Informationen zur Person
beziehungsweise zur sozialen Herkunft. Die Long-Files wurden den klassischen
jahrspezifischen Dateien vorgezogen, weil sie das Stapeln der Wellen überflüssig machen
und damit die Aufbereitung deutlich vereinfachen.

\begin{table}[h]
\centering
\caption{Verwendete Dateien des SHP-Datensatzes}
\label{tab:dateien}
\renewcommand{\arraystretch}{1.2}
\begin{tabularx}{\textwidth}{@{}p{2.8cm} p{2.6cm} X@{}}
\toprule
Datei & Schlüssel & Inhalt \\
\midrule
SHPLONG\_P\_\allowbreak USER.sav &
idpers, year, idhous &
Personenbezogene Variablen je Welle: Vertrauen in den Bundesrat (PP04), politische Selbsteinstufung (PP10), Bildung (EDUCAT), eigene Kinder (OWNKID), Nationalität (NAT\_1\_, NAT\_2\_, NAT\_3\_), gesundheitliche Einschränkung (PC19/PC19A), Zivilstand (CIVSTA). \\
\addlinespace
SHPLONG\_H\_\allowbreak USER.sav &
idhous, year &
Haushaltsbezogene Variablen je Welle: Kanton (CANTON), Netto-Haushaltseinkommen (IHTYN), Haushaltsgrösse (NBPERS). \\
\addlinespace
SHP\_MP.sav &
IDPERS &
Zeitinvariante Masterdaten pro Person: Geburtsjahr (BIRTHY), Geschlecht (SEX). \\
\addlinespace
SHP\_SO.sav &
IDPERS &
Daten zur sozialen Herkunft: finanzielle Probleme in der Jugend (PO58). \\
\bottomrule
\end{tabularx}
\end{table}

Abbildung~\ref{fig:joinschema} fasst schematisch zusammen, über welche
Schlüsselvariablen die vier Dateien zusammengeführt werden. Alle Verknüpfungen
erfolgen als Left Join ausgehend von SHPLONG\_P\_USER, sodass die
Ausgangsmenge der Personen-Jahr-Beobachtungen erhalten bleibt.

\begin{figure}[h]
\centering
\begin{lstlisting}[basicstyle=\ttfamily\small, frame=single, backgroundcolor=\color{white}, numbers=none]
SHPLONG_P_USER  (idpers, year, idhous, pp04, ...)
       |
       |  LEFT JOIN on (idhous, year)
       v
SHPLONG_H_USER  (idhous, year, canton, ihtyn, nbpers)
       |
       |  LEFT JOIN on idpers
       v
SHP_MP          (IDPERS, BIRTHY, SEX)
       |
       |  LEFT JOIN on idpers
       v
SHP_SO          (IDPERS, PO58)
       |
       v
 Filter: year in {1999..2009, 2011, 2014, 2017, 2020, 2023}
 Filter: AGE = year - BIRTHY >= 18
 Konstruktion: BIRTHYGEN, SEX, CIVSTA, CANTON, EDUCAT,
               HASKID, NAT, PC19, NETHHINCPP, PO58, COMPLETE
       |
       v
 panel_final (Personen-Jahr-Ebene)
\end{lstlisting}
\caption{Schema zum Zusammenführen der vier SHP-Dateien zu einem Paneldatensatz.}
\label{fig:joinschema}
\end{figure}

\subsection{Datenkatalog}
\label{subsec:datenkatalog}

Der aufbereitete Datensatz liegt auf Personen-Jahr-Ebene vor. Jede Zeile beschreibt
eine volljährige Person in einem Jahr, in dem die Vertrauensfrage erhoben wurde. Der
Datensatz umfasst 162\,593 Beobachtungen auf 17 Variablen. Ein vollständiger
Datenkatalog mit Wertebereichen, Quelldateien und Rekodierungsregeln findet sich in
Anhang~\ref{AppendixA}. Tabelle~\ref{tab:datenkatalog-kurz} fasst die Variablen
vorab zusammen.

\begin{table}[h]
\centering
\caption{Variablen des aufbereiteten Datensatzes (Kurzfassung)}
\label{tab:datenkatalog-kurz}
\begin{tabularx}{\textwidth}{lX}
\toprule
Variable & Inhalt \\
\midrule
IDPERS     & Personen-ID (Integer). \\
IDHOUS     & Haushalts-ID (Integer). \\
YEAR       & Befragungsjahr (1999--2023, nur Jahre mit PP04). \\
PP04       & Zielvariable: Vertrauen in den Bundesrat, 0--10. \\
BIRTHY     & Geburtsjahr. \\
BIRTHYGEN  & Konstruierte Geburtskohorte in Zehnjahresabschnitten. \\
SEX        & Geschlecht (männlich, weiblich, andere). \\
CANTON     & Wohnkanton (26 Kürzel). \\
EDUCAT     & Höchste abgeschlossene Ausbildung (11 Stufen). \\
HASKID     & Mindestens ein eigenes Kind (ja/nein). \\
NAT        & Nationalitätengruppe (CH, CH+Andere, Andere). \\
PC19       & Langfristige gesundheitliche Einschränkung (ja/nein). \\
PP10       & Politische Selbsteinstufung links--rechts, 0--10. \\
NETHHINCPP & Netto-Haushaltseinkommen pro Kopf (CHF pro Jahr). \\
PO58       & Finanzielle Probleme in der Jugend (ja/nein). \\
COMPLETE   & Kennzeichen vollständiger Beobachtungen (TRUE/FALSE). \\
CIVSTA     & Zivilstand (zusätzliche Einflussvariable). \\
\bottomrule
\end{tabularx}
\end{table}

\subsection{Übersicht der R-Skripte}
\label{subsec:rskripte}

Die Datenverarbeitung ist auf zwei R-Markdown-Skripte aufgeteilt. Das erste Skript
(Datenaufbereitung.Rmd) übernimmt Entpacken, Einlesen, Bereinigung, Merging,
Filtering, Variablenkonstruktion und Speicherung. Das zweite Skript
(Datenqualitaet\_\allowbreak Plausibilisierung\_\allowbreak Auswertung.Rmd) setzt auf dem erzeugten
Datensatz auf und führt die Missing-Analyse, die Imputation, den
BFS-Plausibilitätsvergleich sowie die deskriptive Auswertung durch.

\begin{table}[h]
\centering
\caption{Erstellte R-Skripte}
\label{tab:rskripte}
\begin{tabularx}{\textwidth}{lX}
\toprule
Skript & Inhalt \\
\midrule
Datenaufbereitung.Rmd & Entpacken, Einlesen SPSS, Bereinigung, Merging, Filtering, Variablenkonstruktion, Speicherung. \\[4pt]
\makecell[l]{Datenqualitaet\_ \\ Plausibilisierung\_ \\ Auswertung.Rmd} & Missing-Analyse, Imputation, BFS-Vergleich, deskriptive Auswertung, Hypothesenprüfung. \\
\bottomrule
\end{tabularx}
\end{table}

An R-Paketen kommen tidyverse (Datenmanipulation und Visualisierung),
foreign (SPSS-Import), zoo (LOCF-Imputation), pxR (Einlesen
der BFS-Daten im PX-Format), knitr (Berichterstellung) und scales
(Achsenformatierung) zum Einsatz. Die Auswahl dieser Pakete orientiert sich an den
Empfehlungen der Modulunterlagen.


%----------------------------------------------------------------------------------------

\subsection{Datenaufbereitung: Merging und Filtering}
\label{subsec:merging}

Die Aufbereitung des Paneldatensatzes folgt einer festen Abfolge von Arbeitsschritten,
die vollständig aus R heraus gesteuert werden. Die Reihenfolge entspricht dem in
Abbildung~\ref{fig:joinschema} dargestellten Schema.

Im ersten Schritt wird die SPSS-ZIP-Datei in ein Arbeitsverzeichnis entpackt und die
vier benötigten Dateien mit read.spss() eingelesen. Mit dem Argument
,,use.missings = FALSE`` werden die SHP-spezifischen Missing-Codes, darunter
\mbox{$-1$} (inapplicable), \mbox{$-2$} (no answer), \mbox{$-3$} (does not know) und
\mbox{$-7$} (filter error), automatisch als fehlende Werte markiert. Alle relevanten
Variablen werden anschliessend in numerische Typen umgewandelt, um einheitliche
Datentypen sicherzustellen. Verbliebene negative Codes werden in einem zusätzlichen
Bereinigungsschritt ebenfalls als fehlend gesetzt.

Im zweiten Schritt wird ein Jahrfilter angewendet. Behalten werden ausschliesslich die
16 Wellen, in denen die Zielvariable Vertrauen in den Bundesrat erhoben wurde, nämlich 1999 bis
2009, 2011, 2014, 2017, 2020 und 2023. Die übrigen Wellen enthalten keine
Vertrauensfrage und tragen zur Analyse nichts bei.

Im dritten Schritt werden die vier Dateien schrittweise über Left Joins
zusammengeführt. Das Long-File der Personen wird zuerst mit dem Long-File der
Haushalte über Personen-ID und Befragungsjahr verbunden, sodass jede
Personenbeobachtung die zugehörigen Haushaltsinformationen desselben Jahres erhält.
Anschliessend werden die zeitinvarianten Angaben aus den Masterdaten (Geburtsjahr,
Geschlecht) und aus den Herkunftsdaten (finanzielle Probleme in der Jugend) über die
Personen-ID angehängt. Der Left Join stellt sicher, dass keine
Personenbeobachtung aus dem Personen-Long-File verloren geht.

Im vierten Schritt wird der Erwachsenenfilter gesetzt. Behalten werden nur
Beobachtungen, bei denen das rechnerische Alter mindestens 18 Jahre beträgt. Im fünften
Schritt werden die Zielvariablen gemäss Aufgabenstellung konstruiert. Die Nationalität
wird aus bis zu drei Codes zu den Kategorien CH, CH+Andere und Andere verdichtet, wobei
der Schweizer Code (756 bzw.\ 8100) als Entscheidungskriterium dient. Der Kinderstatus
wird aus der Anzahl eigener Kinder abgeleitet, die Geburtskohorte aus dem Geburtsjahr
als Intervall in Zehnjahresschritten. Das pro Kopf berechnete Netto-Haushaltseinkommen
ergibt sich als Quotient aus dem jährlichen Haushaltseinkommen und der Haushaltsgrösse.
Zum Abschluss wird ein Vollständigkeitskennzeichen gesetzt, das angibt, ob in einer
Zeile alle 15 inhaltlichen Variablen gefüllt sind.

Im sechsten Schritt wird der aufbereitete Datensatz sowohl als CSV- wie auch als
RDS-Datei unter data/processed/ gespeichert. Er umfasst 162\,593
Personenjahre auf 17 Variablen und bildet die gemeinsame Grundlage für alle
nachfolgenden Auswertungen.


%----------------------------------------------------------------------------------------

\subsection{Datenqualität: Fehlende Werte}
\label{subsec:datenqualitaet}

Die Prüfung der Datenqualität konzentriert sich auf fehlende Werte. Tabelle~\ref{tab:missing}
zeigt für jedes Erhebungsjahr die Anzahl der Beobachtungen sowie den Anteil jener Zeilen,
bei denen mindestens eine der 15 inhaltlichen Variablen fehlt.

\begin{table}[h]
\centering
\caption{Anteil Datensätze mit mindestens einem Missing pro Jahr}
\label{tab:missing}
\begin{tabular}{rrrr@{\hspace{2em}}rrrr}
\toprule
Jahr & $n$ & $n_{\text{miss.}}$ & \% &
Jahr & $n$ & $n_{\text{miss.}}$ & \% \\
\midrule
1999 & 9\,633  & 4\,621  & 48.0 & 2005 & 9\,055  & 8\,617  & 95.2 \\
2000 & 8\,719  & 8\,366  & 95.9 & 2006 & 8\,706  & 8\,268  & 95.0 \\
2001 & 8\,168  & 7\,887  & 96.6 & 2007 & 8\,478  & 8\,036  & 94.8 \\
2002 & 8\,019  & 7\,750  & 96.6 & 2008 & 8\,376  & 7\,882  & 94.1 \\
2003 & 7\,837  & 7\,545  & 96.3 & 2009 & 8\,119  & 7\,615  & 93.8 \\
2004 & 7\,593  & 7\,226  & 95.2 & 2011 & 14\,030 & 11\,746 & 83.7 \\
     &         &         &      & 2014 & 14\,943 & 11\,414 & 76.4 \\
     &         &         &      & 2017 & 12\,999 & 9\,139  & 70.3 \\
     &         &         &      & 2020 & 10\,953 & 6\,265  & 57.2 \\
     &         &         &      & 2023 & 10\,965 & 5\,649  & 51.5 \\
\bottomrule
\end{tabular}
\end{table}

Auffällig ist der sprunghafte Anstieg des Missing-Anteils von 48\,\% im Jahr 1999 auf
über 95\,\% ab dem Jahr 2000. Die Auswertung der Missings pro Variable und Jahr zeigt,
dass dieser Anstieg massgeblich auf die Variable langfristige Gesundheitseinschränkung
zurückzuführen ist. Sie wird im SHP nur in den Jahren 1999, 2004, 2009, 2011, 2014,
2017, 2020 und 2023 erhoben. In den übrigen Jahren fehlt die Angabe vollständig, was
nach Konstruktion zu einem Anteil unvollständiger Fälle von nahezu 100\,\% führt.
Ab 2011 fällt der Missing-Anteil allmählich, was mit den Auffrischungsstichproben
(Refreshment Samples SHP III 2013 und SHP IV 2020) und der damit
verbundenen besseren Abdeckung dieser Variable zusammenhängt.

Für die Imputation wird die Variable mit dem höchsten Missing-Anteil im jeweiligen
Zieljahr ausgewählt und mit dem Verfahren Last Observation Carried Forward (LOCF)
behandelt. Für jede Person wird dabei der letzte gültige Wert aus einer früheren
Welle übernommen. Das Verfahren ist konservativ und unterstellt, dass sich der
Zustand zwischen zwei Erhebungen nicht verändert. Nach der Imputation wird das Vollständigkeitskennzeichen neu berechnet, wodurch der
Anteil vollständiger Fälle für das jeweilige Zieljahr steigt.


%----------------------------------------------------------------------------------------

\subsection{Plausibilität: Vergleich mit BFS-Daten}
\label{subsec:plausibilitaet}

Zur Einschätzung der Stichprobenqualität werden die Verteilungen nach Geschlecht,
Altersgruppe und Kanton mit Referenzzahlen des Bundesamts für Statistik verglichen.
Als Referenz dient die ständige Wohnbevölkerung der Schweiz im Jahr 2023 gemäss
STATPOP \cite{BFS2023}. Die BFS-Daten werden über die offene BFS-API direkt aus R im
PX-Format geladen, was die Reproduktion des Vergleichs erleichtert. Verglichen werden
jeweils drei Verteilungen: der gesamte SHP-Datensatz, die Teilmenge der vollständigen
Fälle und die BFS-Referenz.

\paragraph{Geschlecht.}
Im SHP liegen alle Fälle bei 48{,}1\,\% Männern und 51{,}9\,\% Frauen. Die BFS-Referenz
weist 49{,}7\,\% Männer und 50{,}3\,\% Frauen aus. Nach dem Entfernen unvollständiger
Fälle steigt der Frauenanteil auf 52{,}6\,\%. Der Unterschied zur Referenz ist gering
und entspricht dem bekannten, leicht erhöhten Frauenanteil in Haushaltsbefragungen.
Eine systematische Verzerrung ist nicht erkennbar.

\paragraph{Altersgruppen.}
Bei den Altersgruppen sind die Abweichungen deutlicher. Die Gruppe 18 bis 29 Jahre ist
im SHP mit 18{,}1\,\% aller Fälle leicht über, bei den vollständigen Fällen hingegen
deutlich unter der BFS-Referenz von 16{,}5\,\% vertreten (nur 11{,}2\,\%). Die
Altersgruppe 45 bis 59 ist im SHP mit 29{,}2\,\% gegenüber 25{,}7\,\% in der
BFS-Referenz leicht übervertreten. Das Muster entspricht der bekannten
Panelattrition: Jüngere Befragte steigen häufiger aus der Studie aus, während mittlere
und ältere Altersgruppen im Panel bleiben.

\paragraph{Kantone.}
Die Kantonsverteilung im SHP entspricht weitgehend der BFS-Referenz. Die
bevölkerungsreichsten Kantone Zürich, Bern und Waadt sind in beiden Quellen am
stärksten vertreten. Geringe Abweichungen bei kleineren Kantonen lassen sich durch die
geschichtete Stichprobenziehung des SHP erklären, die einzelne Regionen gezielt
übergewichtet, damit auch für kleinere Kantone belastbare Fallzahlen vorliegen.

Insgesamt bildet der SHP-Datensatz die Schweizer Wohnbevölkerung hinsichtlich
Geschlecht und Kanton gut ab. Bei den Altersgruppen ist die für Panels typische
Unterrepräsentation jüngerer und Überrepräsentation mittlerer Personen zu beachten.
Für repräsentative Schätzungen kämen die vom SHP bereitgestellten
Längsschnittgewichte in Frage; für die hier dokumentierte deskriptive Auswertung
werden sie nicht verwendet.
